% id: 170-345
% book: 개념원리 기하 · 21 위치벡터
% section: 개념원리 익히기
% figure: none
% answer: ⑴ $-\vec{a}+3\vec{b}-2\vec{c}$ ⑵ $-5\vec{a}+\vec{b}+4\vec{c}$
% answer_source: 답지
% variant_level: 0 (원본 전사)
% 이 파일은 build.mjs 가 items.json 에서 만든다. 고칠 때는 items.json 을 고친다.
\smash{\raisebox{1.5mm}{\dmkichul{개념원리 기하 170쪽 345번}}}
세 점~$\pt{A}$, $\pt{B}$, $\pt{C}$의 위치벡터를 각각 $\vec{a}$, $\vec{b}$, $\vec{c}$라 할 때, 다음 벡터를 $\vec{a}$, $\vec{b}$, $\vec{c}$로 나타내시오.
\dmsubprob{1} $\vecAB{AB}-2\vecAB{BC}$
\dmsubprob{2} $2\vecAB{AB}+\vecAB{BC}-3\vecAB{CA}$
